\documentclass[11pt]{article}
\usepackage[utf8]{inputenc}
\usepackage{amsthm}


% start: P R E A M B L E ------------------------ P R E A M B L E ------------------------ P R E A M B L E ------------------------

%\input{tikzpreamble.tex}

\title{A Complete Probabilistic Semantics for Lambek Calculus}
\author{}
\date{}

%\usepackage{/Users/lachlanmcpheat/Dropbox/Tex/myStyle}

\renewcommand{\P}{\mathcal{P}}

\usepackage[utf8]{inputenc}
\usepackage[T1]{fontenc}
\usepackage{textcomp}
\usepackage[english]{babel}
\usepackage{amsmath, amssymb}
%\usepackage[left=1in, right=1in]{geometry}
\usepackage{tikz}
%\usepackage{xfp, siunitx}
\usepackage{pgfplots}
\usepackage{proof}
\usepackage{dsfont}

% support for linguistics examples
\usepackage{gb4e}
\noautomath % gb4e allows for the use of _ and ^ outside of math mode by default. \noautomath disables this feature.

% figure support
\usepackage{import}
\usepackage{parskip}
\usepackage{xifthen}
\pdfminorversion=7
\usepackage{pdfpages}
\usepackage{transparent}

\newtheorem{proposition}{Proposition}
\newtheorem{definition}{Definition}
\newtheorem{theorem}{Theorem}[section]
\newtheorem{lemma}{Lemma}


\usepackage{biblatex}


\addbibresource{references.bib}

%\newtheorem*{notation}{Notation}

% end: P R E A M B L E ------------------------ P R E A M B L E ------------------------ P R E A M B L E ------------------------

\begin{document}

\maketitle

\section{Plan}

We extend the capabilities of Lambek Calculus (LC), the logic of syntactic structure used to model natural language. We look into formalising a probabilistic semantics for LC, in the light of probabilistic methods used in natural language processing with great success. We begin by focusing on parsing, the idea of breaking down a sentence into its syntactic structure. The main limitation of parsers based on LC is that a sentence may have ambiguities, resulting in several possible meanings for the sentence. With a probabilistic parser, ambiguities can be resolved by picking the derivation of a sentence with the highest probability.

We start by constructing a formal way of assigning probabilities to words and extending this to sentences (these probabilities can be found by training on a large corpus of text). We then prove the soundness and completeness of LC with respect to these semantics. 

Once this is done, we can compare with existing probabilistic semantics. PCFGs (Probabilistic Context Free Grammars) are a notable example. These are a probabilistic extension of CFGs (Context Free Grammars). The equivalence between the expressive powers of Lambek calculus and CFGs was proven by Pentus \cite{Pentus1993}. We can look into proving a similar equivalence between our probabilistic semantics for LC and PCFGs. 

\section{Formalisation}

We start with the language for LC, which we denote by ${\cal L}_{\text{LC}}$.  The set of atomic types is denoted by $\mathcal{A}$. The logic we use will be Associative LC (might need to tweak a bit in line with Fitting). The axioms are as follows:

\begin{align}
(a)\quad & x \to x  \\[4pt]
(b)\quad & (xy)z \to x(yz)
& (b')\quad & x(yz) \to (xy)z \\[4pt]
(c)\quad & \text{if } xy \to z \text{ then } x \to z/y
& (c')\quad & \text{if } xy \to z \text{ then } y \to x\backslash z \\[4pt]
(d)\quad & \text{if } x \to z/y \text{ then } xy \to z
& (d')\quad & \text{if } y \to x\backslash z \text{ then } xy \to z \\[4pt]
(e)\quad & {\text{if } x \to y \text{ and } y \to z \text{ then } x \to z}
\end{align}

We now work towards constructing a probablistic semantics for LC. We start with a corpus $\mathcal{C}$ of sentences. Let $\Sigma$ be the vocabulary consisting of all words found in $\mathcal{C}$. We are interested in parsing phrases or complete sentences, i.e  elements of $\Sigma^*$, which is the set of all strings from $\Sigma$. Note that $\Sigma^*$ is the superset of ${\cal C}$. 

A frame over $\Sigma$ is the structure $\langle \Sigma^\ast, R \rangle$ where $R$ is a ternary relation on $\Sigma^\ast$. Intuitively, $R(x, y, z)$ means that string $z$ can be decomposed as $x \cdot y$ (i.e., as $x$ followed by $y$). Note that, in general, there are multiple ways of decomposing a string in this way. For example, the following both hold:

\[R(\textbf{`cats'}, \textbf{ `and dogs' }, \textbf{ `cats and dogs'})\]
\[R(\textbf{`cats and'}, \textbf{ `dogs' }, \textbf{ `cats and dogs'})\]

A valuation is a map $v \colon \Sigma^\ast \times \mathcal{T} \to [0,1]$ giving
the probability $v(x, A)$ of a string $x$ having a parse structure, or type,
$A$, relative to our corpus $\mathcal{C}$. If $x \in \Sigma^\ast$ is an atomic
string (a single word) then $v(x, A)$ is computed from the frequency of
occurrences of $x$ in $\mathcal{C}$.  Let $\mathcal{T}_{\mathcal{C}} (x)$ denote
the set of all types  the word $x$ has in the corpus $\mathcal{C}$; we,
naturally, require that
\begin{equation}
  \sum_{A \in \mathcal{T}_{\mathcal{C}} (x)} v(x,A) = 1.
\end{equation}
For words $x' \in {\cal C}$ that are in the corpus but never have a type $A'$, the probability should be zero, i.e. $v(x', A') = 0$. If the word never occurred in the corpus,  i.e. $x'' \notin {\cal C}$, its probability is undefined and we assign the value $+\infty$ to it, i.e. $v(x'', X) = +\infty$, for all types $X \in {\cal L}_{\text{LC}}$. 

For divisible strings, the valuation is computed recursively:
\begin{itemize}
    \item $v(x, A \circ B) := \max \{ v(y, A) \cdot v(z, B) \mid R(y, z, x), \ \text{where} \ y, z \in \Sigma^* \}$;
    \item $v(x, A \setminus B) := \min \{ \frac{v(z, B)}{v(y, A)} \mid R(y, x, z),
      \mbox{with $z \in \mathcal{C}$}, v(y, A) \neq 0 \}$ (we restrict ourselves to $z \in
      \mathcal{C}$ so that we do not have to consider infinitely many strings $z
      \in \Sigma^\ast$) \\ Note also that this definition makes sense if we have the following:
      \[\exists z \in \mathcal{C},  y \in \Sigma^* \; s.t \; R(y, x, z) , v(y, A) > 0, \frac{v(z, B)}{v(y, A)} \leq 1\]
    \item $v(x, B / A) := \min \{ \frac{v(z, B)}{v(y, A)} \mid R(x, y, z),
      \mbox{with $z \in \mathcal{C}$} \}$ (note that the order of the arguments
      of the relation $R$ differs from that in the preceding clause).
\end{itemize}


Note that if $v(y,A) = 0$, then $v(z, B) \div v(y, A)$, being infinity, does not
affect the value of $\min \{ v(z, B) \div v(y, A) \mid R(x, y, z), z \in C \}$;
however, for this value to be bounded above by $1$,  we need, for every $z$, at
least one $y$ with  $v(y,A) \ne 0$.

We say that $A$ implies $B$, and write $A \models B$, if $v(x, A) \leqslant
v(x,B)$, for every $x \in \Sigma^\ast$. {\color{red} Before we considered all strings, but because of complications when trying to prove soundness, we restrict to strings consisting of words in corpus}

RESTTRICT TO CORPUS

We next show that LC is sound and complete with respect to this semantics.

\section{Soundness}

We want to show that $A \vdash_{LC} B \Rightarrow A \models B$. It is sufficent to show that all the axioms (a)-(e) for LC are valid in our semantics, as all other statements in LC are derived from these axioms. 

(a) $A \models A$
\begin{proof}
Let $x \in \Sigma^\ast$. Obviously, $v(x, A) \leq v(x, A)$ so we are done.
\end{proof}

(b) $(A \circ B) \circ C \models A \circ (B \circ C)$
\begin{proof}
Let $x \in \Sigma^\ast$. 

\[ v(x, (A \circ B) \circ C))\]
\[ = \max \{ v(y, A \circ B) \cdot v(c, C) \mid R(y, c, x) \}\]
\[ = \max \{ \max \{ v(a, A) \cdot v(b, B) \mid R(a, b, y) \} \cdot v(c, C) \mid R(y, c, x) \}\]
\[ = \max \{ v(a, A) \cdot v(b, B)  \cdot v(c, C) \mid R(y, c, x), \; R(a, b, y)\}\]
\[ = \max \{ v(a, A) \cdot v(b, B)  \cdot v(c, C) \mid R(a, z, x), \; R(b, c, z)\}  \tag{
$\ast$}\]
\[ = \max\{v(a, A) \cdot  \max \{ v(b, B) \cdot v(c, C) \mid R(b, c, z) \} \mid R(a, z, x) \}\]
\[ = \max \{ v(a, A) \cdot v(z, B \circ C) \mid R(a, z, x) \}\]
\[ = v(x, A \circ (B \circ C))\]

{\color{red}TODO explain the ($\ast$) and double check the nested max replacement}

\end{proof}

(c), (d) $A \circ B \models C \iff A \models C/B$
\begin{proof}

$\forall x\; v(x, A \circ B) \leq v(x, C)$
\[\iff \forall x\; \max \{ v(y, A) \cdot v(z, B) \mid R(y, z, x) \} \leq v(x, C)\]
\[\iff \forall x, y, z \; s.t \; R(y, z, x), \; v(y, A) \cdot v(z, B) \leq v(x, C)\]
\[\iff \forall x, y, z \; s.t \; R(y, z, x), \; v(y, A) \leq \frac{v(x, C)}{v(z, B)}\]
\[\iff \forall y \; v(y, A) \leq \min \{ \frac{v(x, C)}{v(z, B)} \mid R(y, z, x) \}\]
\[\iff \forall y \; v(y, A) \leq \min \{ \frac{v(x, C)}{v(z, B)} \mid R(y, z, x)\} = v(y, C/B) \tag{$\ast$ $\ast$} \] 
\[\iff A \models C/B\]
As required.
\end{proof}

($\ast$) {\color{red}We assume that division by zero yields $+\infty$}

($\ast$ $\ast$) {\color{red}This skips the subtlety that $x$ must be in the corpus $\mathcal{C}$}


(e) $\text{if } A \models B \text{ and } B \models C\text{ then } A \models C$
\begin{proof}
Let $x \in \Sigma^\ast$.
If $v(x, A) \leq v(x, B)$ and $v(x, B) \leq v(x, C)$ then by transitivity we have $v(x, A) \leq v(x, C)$ and the statement follows.
\end{proof}



\end{document}
%%% Local Variables:
%%% mode: latex
%%% TeX-master: t
%%% End:
